% GENERATED FILE. Edit canonical lessons, then run npm run generate:books.
% canonical-source-hash: fnv1a64:ee3d98638529fa45
% canonical-lessons: MR-C300-dhamkavne, MR-C300-lac, MR-C300-siksa, MR-C300-pal-kadhne, MR-C300-kondne

\chapter{Threats and Bribes}
\label{ch:mr-threats-and-bribes}
\hlchaptermodality{300}

\begin{chapteropening}
\textbf{By the end of this chapter:} \emph{I can talk about threats, bribes and punishment.}

Complete the last lesson of chapter 300: I can talk about threats, bribes and punishment.
\end{chapteropening}

\section[\texorpdfstring{\mr{धमकावणे} (dhamkāvṇe)}{धमकावणे (dhamkāvṇe)}]{\mr{धमकावणे} (dhamkāvṇe) — to threaten}

\label{lesson:MR-C300-dhamkavne}



\begin{quote}
Before the new one: say the Marathi for to confess, then the Marathi for to file.
\end{quote}

\subsection*{You'll want to know: \mr{धमकावणे}}
\textbf{\mr{धमकावणे}} — \emph{dhamkāvṇe} — \textquotedblleft{}to threaten\textquotedblright{}.

\begin{quote}
\mr{त्या} \mr{गुंडाने} \mr{दुकानदाराला} \mr{धमकावलं}.
\end{quote}

\begin{quote}
\emph{tyā guṇḍāne dukāndārālā dhamkāvlaṁ.}
\end{quote}

\begin{quote}
That thug threatened the shopkeeper.
\end{quote}

Past, with -\mr{ने} (he …): \mr{त्याने} \mr{धमकावलं}, \emph{tyāne dhamkāvlaṁ}.

The person threatened takes -\mr{ला}; the noun is \mr{धमकी}, a threat.

\begin{cousinweb}
\textbf{\mr{धमकावणे}} (\emph{dhamkāvṇe}) is inherited: CDIAL lists Marathi \emph{dhamkāvṇe}, 'to threaten', under \emph{*dhammakka}, 'a threat' (CDIAL 6736).
\end{cousinweb}

\begin{grammarlens}[title={What you've built}]
One more word for telling what happened.
\end{grammarlens}

\subsection*{Your turn}
\begin{itemize}
\raggedright
  \item \emph{Say it:} \emph{dhamkāvṇe}
  \item \emph{Say it:} \emph{dhamkāvṇe}, once more
  \item \emph{Say it:} \emph{tyā guṇḍāne dukāndārālā dhamkāvlaṁ.} — That thug threatened the shopkeeper.
  \item \emph{Recall:} say the Marathi for to confess, then the Marathi for to file, then say \emph{dhamkāvṇe} again
\end{itemize}

\begin{tcolorbox}[breakable,colback=teal!4,colframe=teal!35!black,title={Before you move on}]
What does \mr{धमकावणे} mean? (\textquotedblleft{}To threaten\textquotedblright{}.) Say it once more.
\end{tcolorbox}
\section[\texorpdfstring{\mr{लाच} (lāc)}{लाच (lāc)}]{\mr{लाच} (lāc) — a bribe}

\label{lesson:MR-C300-lac}



\begin{quote}
Before the new one: say the Marathi for to file, then the Marathi for to threaten.
\end{quote}

\subsection*{You'll want to know: \mr{लाच}}
\textbf{\mr{लाच}} — \emph{lāc} — \textquotedblleft{}a bribe\textquotedblright{}.

\begin{quote}
\mr{लाच} \mr{घेताना} \mr{एका} \mr{कारकुनाला} \mr{पकडलं}.
\end{quote}

\begin{quote}
\emph{lāc ghetānā ekā kārkunālā pakaḍlaṁ.}
\end{quote}

\begin{quote}
A clerk was caught taking a bribe.
\end{quote}

Feminine.

\mr{लाच} \mr{घेणे}, to take a bribe; \mr{लाच} \mr{देणे}, to give one.

\begin{cousinweb}
\textbf{\mr{लाच}} (\emph{lāc}) is a tadbhava, inherited from Sanskrit \emph{lañcā}, 'a present, a bribe' (CDIAL 10914).
\end{cousinweb}

\begin{grammarlens}[title={What you've built}]
One more word for telling what happened.
\end{grammarlens}

\subsection*{Your turn}
\begin{itemize}
\raggedright
  \item \emph{Say it:} \emph{lāc}
  \item \emph{Say it:} \emph{lāc}, once more
  \item \emph{Say it:} \emph{lāc ghetānā ekā kārkunālā pakaḍlaṁ.} — A clerk was caught taking a bribe.
  \item \emph{Recall:} say the Marathi for to file, then the Marathi for to threaten, then say \emph{lāc} again
\end{itemize}

\begin{tcolorbox}[breakable,colback=teal!4,colframe=teal!35!black,title={Before you move on}]
What does \mr{लाच} mean? (\textquotedblleft{}A bribe\textquotedblright{}.) Say it once more.
\end{tcolorbox}
\section[\texorpdfstring{\mr{शिक्षा} (śikṣā)}{शिक्षा (śikṣā)}]{\mr{शिक्षा} (śikṣā) — a punishment; a sentence}

\label{lesson:MR-C300-siksa}



\begin{quote}
Before the new one: say the Marathi for to threaten, then the Marathi for a bribe.
\end{quote}

\subsection*{You'll want to know: \mr{शिक्षा}}
\textbf{\mr{शिक्षा}} — \emph{śikṣā} — \textquotedblleft{}a punishment; a sentence\textquotedblright{}.

\begin{quote}
\mr{न्यायालयाने} \mr{त्याला} \mr{तीन} \mr{वर्षांची} \mr{शिक्षा} \mr{दिली}.
\end{quote}

\begin{quote}
\emph{nyāyālayāne tyālā tīn varṣāñcī śikṣā dilī.}
\end{quote}

\begin{quote}
The court sentenced him to three years.
\end{quote}

Feminine.

\mr{शिक्षा} \mr{देणे}, to sentence; \mr{शिक्षा} \mr{होणे}, to be punished. In Marathi it means punishment, not teaching.

\begin{cousinweb}
\textbf{\mr{शिक्षा}} (\emph{śikṣā}) is a tatsama, taken from Sanskrit \emph{śikṣā}, 'learning, teaching; also chastisement, punishment' (Monier-Williams); Marathi keeps the last sense.
\end{cousinweb}

\begin{grammarlens}[title={What you've built}]
One more word for telling what happened.
\end{grammarlens}

\subsection*{Your turn}
\begin{itemize}
\raggedright
  \item \emph{Say it:} \emph{śikṣā}
  \item \emph{Say it:} \emph{śikṣā}, once more
  \item \emph{Say it:} \emph{nyāyālayāne tyālā tīn varṣāñcī śikṣā dilī.} — The court sentenced him to three years.
  \item \emph{Recall:} say the Marathi for to threaten, then the Marathi for a bribe, then say \emph{śikṣā} again
\end{itemize}

\begin{tcolorbox}[breakable,colback=teal!4,colframe=teal!35!black,title={Before you move on}]
What does \mr{शिक्षा} mean? (\textquotedblleft{}A punishment; a sentence\textquotedblright{}.) Say it once more.
\end{tcolorbox}
\section[\texorpdfstring{\mr{पळ} \mr{काढणे} (paḷ kāḍhṇe)}{पळ काढणे (paḷ kāḍhṇe)}]{\mr{पळ} \mr{काढणे} (paḷ kāḍhṇe) — to flee; to take to one's heels}

\label{lesson:MR-C300-pal-kadhne}



\begin{quote}
Before the new one: say the Marathi for a bribe, then the Marathi for a punishment.
\end{quote}

\subsection*{You'll want to know: \mr{पळ} \mr{काढणे}}
\textbf{\mr{पळ} \mr{काढणे}} — \emph{paḷ kāḍhṇe} — \textquotedblleft{}to flee; to take to one's heels\textquotedblright{}.

\begin{quote}
\mr{पोलिसांना} \mr{पाहिल्यावर} \mr{चोराने} \mr{पळ} \mr{काढला}.
\end{quote}

\begin{quote}
\emph{polisānnā pāhilyāvar corāne paḷ kāḍhlā.}
\end{quote}

\begin{quote}
When he saw the police, the thief fled.
\end{quote}

Past, with -\mr{ने} (he …): \mr{त्याने} \mr{पळ} \mr{काढला}, \emph{tyāne paḷ kāḍhlā}.

The agent takes -\mr{ने} and \mr{पळ} is masculine: \mr{चोराने} \mr{पळ} \mr{काढला}. Plainer: \mr{पळून} \mr{जाणे}.

\begin{cousinweb}
\textbf{\mr{पळ} \mr{काढणे}} (\emph{paḷ kāḍhṇe}) is built from \emph{paḷ}, 'flight', of the family of Sanskrit \emph{palāyate}, 'flees' (CDIAL 7955), and \emph{kāḍhṇe}, 'to take out'.
\end{cousinweb}

\begin{grammarlens}[title={What you've built}]
One more word for telling what happened.
\end{grammarlens}

\subsection*{Your turn}
\begin{itemize}
\raggedright
  \item \emph{Say it:} \emph{paḷ kāḍhṇe}
  \item \emph{Say it:} \emph{paḷ kāḍhṇe}, once more
  \item \emph{Say it:} \emph{polisānnā pāhilyāvar corāne paḷ kāḍhlā.} — When he saw the police, the thief fled.
  \item \emph{Recall:} say the Marathi for a bribe, then the Marathi for a punishment, then say \emph{paḷ kāḍhṇe} again
\end{itemize}

\begin{tcolorbox}[breakable,colback=teal!4,colframe=teal!35!black,title={Before you move on}]
What does \mr{पळ} \mr{काढणे} mean? (\textquotedblleft{}To flee; to take to one's heels\textquotedblright{}.) Say it once more.
\end{tcolorbox}
\section[\texorpdfstring{\mr{कोंडणे} (koṇḍṇe)}{कोंडणे (koṇḍṇe)}]{\mr{कोंडणे} (koṇḍṇe) — to shut in; to lock up}

\label{lesson:MR-C300-kondne}



\begin{quote}
Before the new one: say the Marathi for a punishment, then the Marathi for to flee.
\end{quote}

\subsection*{You'll want to know: \mr{कोंडणे}}
\textbf{\mr{कोंडणे}} — \emph{koṇḍṇe} — \textquotedblleft{}to shut in; to lock up\textquotedblright{}.

\begin{quote}
\mr{चोरांनी} \mr{चौकीदाराला} \mr{खोलीत} \mr{कोंडलं}.
\end{quote}

\begin{quote}
\emph{corānnī caukīdārālā kholīt koṇḍlaṁ.}
\end{quote}

\begin{quote}
The thieves locked the watchman in a room.
\end{quote}

Past, with -\mr{ने} (he …): \mr{त्याने} \mr{कोंडलं}, \emph{tyāne koṇḍlaṁ}.

The place takes -\mr{त}: \mr{खोलीत} \mr{कोंडणे}; \mr{कोंडून} \mr{ठेवणे}, to keep shut in.

\begin{cousinweb}
\textbf{\mr{कोंडणे}} (\emph{koṇḍṇe}) is a Marathi verb of shutting in; its origin is not settled.
\end{cousinweb}

\begin{grammarlens}[title={What you've built}]
One more word for telling what happened.
\end{grammarlens}

\subsection*{Your turn}
\begin{itemize}
\raggedright
  \item \emph{Say it:} \emph{koṇḍṇe}
  \item \emph{Say it:} \emph{koṇḍṇe}, once more
  \item \emph{Say it:} \emph{corānnī caukīdārālā kholīt koṇḍlaṁ.} — The thieves locked the watchman in a room.
  \item \emph{Recall:} say the Marathi for a punishment, then the Marathi for to flee, then say \emph{koṇḍṇe} again
\end{itemize}

\begin{tcolorbox}[breakable,colback=teal!4,colframe=teal!35!black,title={Before you move on}]
What does \mr{कोंडणे} mean? (\textquotedblleft{}To shut in; to lock up\textquotedblright{}.) Say it once more.
\end{tcolorbox}
